\section{A joint covariance bound for cached contraction}
\label{sec:joint-objective}

The strongest cached incoming certificate has time coefficient
$35.3545880584\ldots$ and space coefficient $34.2748030050\ldots$ on
$s/\sqrt{\log_2s}$.  Its binary table construction and fast triangle
contraction are both useful, but their separate worst-case entropy bounds
need not occur on the same group partition.  The new algorithm chooses
one partition by the largest of \emph{all} its construction and triangle
costs.  A covariance calculation controls this joint objective and pays
for every intermediate table and every recursively evaluated child.

We use the exact pair-index identity~\eqref{eq:mode-graph} and
Lemma~\ref{lem:norm}; write $K(M,u)=\K(M,u)$ in this section.
Every group function is evaluated by a paid recursive child contraction.
The uniform bound needed for partition selection is
\begin{equation}\label{eq:joint-entropy-norm}
 \|h\|_2\le(1/r+o(1))K(M,u).
\end{equation}
The source of the cached interface and exact quotient-ring matrix
bridge is the incoming contraction-tree packet~\cite{incomingtree}.

\subsection{Six costs that must be controlled together}
Let $r=3k$ and equipartition the groups into $A,B,C$, each of size $k$.
Write $I_A,I_B,I_C$ for the total directed-pair entropy incident to each
block, counting an internal pair once.  Write $H_{AB},H_{AC},H_{BC}$
for the three crossing entropies, and put
$S=H_{AB}+H_{AC}+H_{BC}$.  Define
\begin{equation}\label{eq:joint-costs}
 \mathcal C(\mathcal P)=
 \bigl(I_A,I_B,I_C,
       S-\tfrac35H_{AB},S-\tfrac35H_{AC},S-\tfrac35H_{BC}\bigr),
 \quad W(\mathcal P)=\max_{1\le j\le6}\mathcal C_j(\mathcal P).
\end{equation}
The first three costs permit direct construction of the three block
tables from cached group functions.  The largest of the last three is
$S-(3/5)\min(H_{AB},H_{AC},H_{BC})$, the rectangular matrix multiplication
cost supplied by an exponent-$12/5$ backend.  The integer-coordinate/CRT bridge in the incoming report~\cite{incomingtree}
makes that backend valid over the actual power-of-two quotient ring.
Here are its bit-cost and algebraic requirements. The ring is a free module
of polynomial rank over $\mathbb Z/2^p\mathbb Z$, with a polynomial-size
integer multiplication table. Lift coordinates to integers, compute each
coordinate-pair matrix product modulo sufficiently many nonexceptional odd
primes, reconstruct by CRT, apply the structure constants, and reduce
modulo $2^p$. An unreduced entry has bit length $O(p+\log d)$ for inner
dimension $d$, which is polynomial in $s$. A fixed rational bilinear
algorithm of exponent below $12/5$ exists by the established field bound
\cite{mm}; excluding its finitely many denominator primes makes its
identities valid over these fields. Coordinate arithmetic and CRT add only
polynomial bit overhead. Sequential recursion uses a polynomial factor
times the input/output table sizes in storage. The algorithm never inverts
two, a nilpotent, or an arbitrary nonunit in the source ring.

\begin{lemma}[Joint adaptive partition bound]\label{lem:joint-partition}
For $r=9$ or $r=12$, an equipartition can be selected so that
\[
 W(\mathcal P)\le r\beta_r\|h\|_2,
 \qquad
 \beta_9^2=\frac{2489356}{6638625},\qquad
 \beta_{12}^2=\frac{4786024}{13540725}.
\]
In particular $\beta_9=0.6123566464\ldots$ and
$\beta_{12}=0.5945200108\ldots$.
\end{lemma}
\begin{proof}
Combine opposite directed weights into undirected weights
$x_{\{g,h\}}=h_{gh}+h_{hg}$.  Put
$H=\sum_eh_e=\sum_ex_e$, $N=r(r-1)$, and $m=N/2$.
Choose a uniformly random labelled equipartition.  For one undirected
edge its vector of coefficients in~\eqref{eq:joint-costs} is, up to
permuting block labels,
\[
 v_{\rm internal,A}=(1,0,0,0,0,0),\qquad
 v_{AB}=(1,1,0,2/5,1,1).
\]
Center each such six-vector by subtracting its coordinate mean; call
the result $\widetilde v_e$.  If
$\overline{\mathcal C}=\frac16\sum_j\mathcal C_j$ and
$V=\sum_j(\mathcal C_j-\overline{\mathcal C})^2$, then
\[
 \mathbb E\overline{\mathcal C}=A_rH,
 \quad A_r=\frac{1+(17/5)p_r}{6},
 \quad p_r=\frac{2k}{r-1},
 \qquad
 \mathbb EV=\sum_{e,f}x_ex_f\,
              \mathbb E\langle\widetilde v_e,\widetilde v_f\rangle.
\]
The covariance entry depends only on whether $e=f$, whether the edges
share one endpoint, or whether they are disjoint.  Denote the three
entries by $c_0,c_1,c_2$.  They are the following exact rational numbers:
\[
\begin{array}{c|ccc|c}
 r&c_0&c_1&c_2&A_r\\\hline
 9&109/120&869/4200&-379/2100&71/120\\
 12&299/330&181/750&-3151/24750&191/330
\end{array}
\]
For completeness, these entries require only a fixed finite color
calculation.  Take edges $(1,2)$ and respectively $(1,2)$, $(1,3)$,
$(3,4)$.  For their $v=2,3,4$ distinct endpoints, sum over the $3^v$
colorings $\chi$ with colors $A,B,C$.  If $n_j$ endpoints receive color
$j$, the probability is
\[
 \frac{(k)_{n_A}(k)_{n_B}(k)_{n_C}}{(3k)_v},
 \qquad (a)_t=a(a-1)\cdots(a-t+1).
\]
Multiply this by the dot product of the two centered coefficient
vectors and sum.  This evaluates the table using at most $81$ terms per
entry.  For example $c_0=(1-p_r)(5/6)+p_r(14/15)$.
The accompanying checker evaluates every term as a rational number.

Let $x'_e=x_e-H/m$ and $d'_v=\sum_{e\ni v}x'_e$.
The covariance matrix has constant row sum
\[
 L_0=c_0+2(r-2)c_1+\tfrac12(r-2)(r-3)c_2.
\]
On the centered weights its quadratic form is
$(c_0-2c_1+c_2)\|x'\|_2^2+(c_1-c_2)\|d'\|_2^2$.
Cauchy--Schwarz gives
$\|d'\|_2^2\le(r-2)\|x'\|_2^2$: use
$\sum_v(d'_v)^2=\sum_{\{v,w\}}x'_{vw}(d'_v+d'_w)$ and
$\sum_vd'_v=0$.  Since $c_1-c_2\ge0$,
\[
 \mathbb EV\le L_1\left(\|x\|_2^2-\frac{H^2}{m}\right)
                         +L_0\frac{H^2}{m},
 \quad L_1=c_0+(r-4)c_1-(r-3)c_2.
\]
The values are
\[
 (L_0,L_1)=
 \begin{cases}
 (3/200,1059/350),&r=9,\\
 (1/275,5476/1375),&r=12.
 \end{cases}
\]
In both cases $0\le L_0<L_1$.
Because $\|x\|_2^2\le2\|h\|_2^2$, it follows that
\[
 \mathbb EV\le2L_1\|h\|_2^2
                         -\frac{2(L_1-L_0)}{N}H^2.
\]
For six numbers, the largest is at most their mean plus
$\sqrt{5/6}$ times their centered Euclidean norm.  Jensen therefore
implies, with $q=1-L_0/L_1>0$,
\[
 \mathbb EW\le A_rH+
   \sqrt{5L_1/3}\sqrt{\|h\|_2^2-qH^2/N}
 \le\sqrt{\frac{A_r^2N}{q}+\frac53L_1}\,\|h\|_2.
\]
The last step is Cauchy--Schwarz with coordinates
$\sqrt q\,H/\sqrt N$ and
$\sqrt{\|h\|_2^2-qH^2/N}$.
A minimum-cost partition is no worse than the mean.  Substitution of
the rational table gives
$\beta_r^2=r^{-2}(A_r^2N/q+5L_1/3)$, and the displayed fractions.
\end{proof}

\paragraph{Exact partition selection.}
All objectives can be compared without logarithms.  For each block let
$D_I$ be the product of the exact list sizes incident to it, and let
$D_{AB},D_{AC},D_{BC}$ be the three crossing dimensions.  Minimize the
maximum of the six rational scores
\[
 \begin{gathered}
 D_{I_A}^5,\quad D_{I_B}^5,\quad D_{I_C}^5,\\[3pt]
 \frac{(D_{AB}D_{AC}D_{BC})^5}{D_{AB}^3},\quad
 \frac{(D_{AB}D_{AC}D_{BC})^5}{D_{AC}^3},\quad
 \frac{(D_{AB}D_{AC}D_{BC})^5}{D_{BC}^3}.
 \end{gathered}
\]
These scores have logarithms $5\mathcal C_j$.  There are only $280$
unlabelled equipartitions for $r=9$, and $5775$ for $r=12$.  Their counts
and their exact integer bit lengths are polynomially bounded in $s$.

\subsection{Every cache, table, and recursive call is paid}
First cache all group functions $Z_g$ on their incident indices.  The
number of child calls has exponent at most
$(\alpha_r+o(1))K(M,u)$, where
$\alpha_r=\sqrt{2(r-1)}/r$, by~\eqref{eq:joint-entropy-norm}.
Each child has at most $(2/r+o(1))M$ wires.

After choosing the partition, form a block table by summing the product
of its cached group functions over its internal pair indices.  Literal
enumeration costs at most a polynomial factor times $2^{I_B}$.
A sequential implementation expands each factor into records on all
incident indices of the block, sorts on the common assignment key,
multiplies the matched records, then sorts and sums on the boundary
key.  Each expanded list has at most $2^{I_B}$ records.  Thus the
claimed bit work and space do not rely on free random access to
exponential tables.  Boundary indices are a subset of incident indices,
so the output table also has at most $2^{I_B}$ entries.

Contract the three tables by the imported rectangular matrix backend.
Its time exponent is the largest of the three triangle costs in
\eqref{eq:joint-costs}; its input/output dimensions and its sequential
workspace are bounded by the same exponent.  Every group cache is no
larger than the incident-index set of the block containing that group.
The selected $W$ consequently bounds \emph{all} local storage as well
as all nonrecursive work.  Equations~\eqref{eq:joint-entropy-norm} and
Lemma~\ref{lem:joint-partition} prove the local recurrence
\begin{equation}\label{eq:joint-recurrence}
 T_d\le\operatorname{poly}(s)\left(
  2^{(\alpha_{r_d}+o(1))K_d}T_{d+1}
       +2^{(\beta_{r_d}+o(1))K_d}\right),
 \qquad S_{\rm local}\le\operatorname{poly}(s)2^{(\beta_{r_d}+o(1))K_d}.
\end{equation}
For the small late alphabets $r=6,4$ below, a simpler cached algorithm
suffices: cache all group functions and sum over all pair indices.
It has the same child-call coefficient $\alpha_r$, local time
coefficient $\beta_r^{\rm full}=\sqrt{r(r-1)}/r$, and local storage
coefficient at most $\alpha_r$.

\subsection{A multistage schedule with strict packing margins}
Put $\lambda=\lceil\log_2s\rceil$,
$g=\lceil\sqrt\lambda\rceil$, and
$b=\lceil6\log_2\lambda\rceil$.  Let $B$ be the largest nonnegative
integer with $9^{20B}\le2^{\lambda-g}$, using direct contraction for
bounded exceptional sizes.  Define
\[
\begin{array}{c|c|c|c}
 d&r_d&u_d&\text{upper cumulative retained range}\\\hline
 0&9&B&(\lambda-g)/(20\log_2 9)\\
 1&12&\lfloor\lambda/36\rfloor-B&\lambda/36\\
 2&9&\lfloor\lambda/288\rfloor&\lambda/32\\
 3&6&\lfloor5\lambda/864\rfloor&\lambda/27\\
 d\ge4&4&\lfloor\lambda/(189\cdot2^{d-4})\rfloor&\lambda/21
\end{array}
\]
Set $a_d=\sum_{j\le d}(u_j+b)$ with $a_{-1}=0$, and leaf threshold
$H=\lceil s/\lambda^{3/4}\rceil$.
The cumulative caps use
$1/36+1/288=1/32$,
$1/32+5/864=1/27$, and $1/27+2/189=1/21$.

At the root the symmetric term $Q_0^2\Delta$ has binary logarithm at
most $\lambda-g/10+O(\log\lambda)$.  At later depths the retained-side
condition has exponent $9/10+(\log_2r_d)a_d/\lambda$.
Strict rational margins follow from
\[
 12^{100}<2^{359},\quad9^{40}<2^{127},\quad6^5<2^{13},
\]
which give respectively margins $1/3600$, $1/1280$, and $1/270$ at
depths $1,2,3$; the constant-$4$ tail has margin $1/210$.
The fresh-side and padding exponents are at most
$(2\log_2r_d+1)a_d/\lambda<1$ throughout.  Fresh packet degrees have
exponent at most $1/21+o(1)<9/10$.

For a noncircular bootstrap, consider hypothetical active depths through
$4+\tfrac34\log_{9/5}\lambda+O(1)$.
The tail ranges satisfy
\[
 u_d\ge\lambda^{\eta+o(1)},\qquad
 \eta=1-\frac{3\ln2}{4\ln(9/5)}>0,
\]
since $2^3<(9/5)^4$.  All packing gates hold uniformly on this range.
Nearly lossless padding then gives eventual tail shrinkage at most
$5/9$, closing the hypothesized depth bound.  The sharp masses on a
path are, up to uniformly vanishing multiplicative perturbations,
\[
 m_0=1,\quad m_1=2/9,\quad m_2=1/27,\quad
 m_3=2/243,\quad m_{4+t}=\frac{2}{729\,2^t},
 \qquad M_d\le(1+o(1))^d s m_d.
\]
The additive buffer gives $D\ge2^{a_{d-1}+b+1}$ at retained levels.
The inherited precision estimate $O(s\lambda^2 2^{a_{d-1}})$ is
$o(MD/u)$ because $s/M\le\lambda^{3/4}$ and
$\lambda^{15/4}2^{-b}\to0$.  Thus the entropy bounds used above are
uniform at every active call.

\subsection{Certified total time and space}
Let $\kappa=\pi\sqrt2/\ln2$ and $A=1/(20\log_2 9)$.  On the scale
$s/\sqrt\lambda$, the $K_d$ coefficients are
\[
 K_0^*=\kappa/\sqrt A,\quad
 K_1^*=\frac{2\kappa}{9\sqrt{1/36-A}},\quad
 K_2^*=\frac{\kappa\sqrt{288}}{27},\quad
 K_3^*=\frac{2\kappa}{243}\sqrt{864/5},\quad
 K_{4+t}^*=\frac{2\kappa\sqrt{189}}{729}\,2^{-t/2}.
\]
Write $P_d=\sum_{j<d}\alpha_{r_j}K_j^*$ for $0\le d\le4$.
Expanding~\eqref{eq:joint-recurrence} gives the following complete list
of competing coefficients:
\[
\begin{array}{c|c}
 \text{source of work}&\text{coefficient}\\\hline
 \text{root local work}&\beta_9K_0^*=31.25239731265\ldots\\
 \text{depth 1 local work}&P_1+\beta_{12}K_1^*=30.41173930826\ldots\\
 \text{depth 2 local work}&P_2+\beta_9K_2^*=30.23123727305\ldots\\
 \text{depth 3 full pair sum}&P_3+(\sqrt{30}/6)K_3^*=30.18781782003\ldots\\
 \text{entire constant-4 tail}&
 P_4+K_4^*\max\{\sqrt{12}/4,(\sqrt6/4)/(1-1/\sqrt2)\}
 =30.42570652863\ldots
\end{array}
\]
Exact rational interval arithmetic in the checker verifies every strict
comparison.  Hence
\begin{equation}\label{eq:joint-headline}
 C_{\rm joint}=
 \frac{\pi\sqrt2}{\ln2}\sqrt{20\log_2 9}
       \sqrt{\frac{2489356}{6638625}}
 =31.2523973126509\ldots
\end{equation}
bounds the total time coefficient.  Local storage is bounded by
$\beta_9K_0^*$, $\beta_{12}K_1^*$, $\beta_9K_2^*$, and the smaller
late singleton caches; ancestor tables add rather than multiply.
The same $C_{\rm joint}$ is therefore a valid space coefficient.
This space bound is an upper certificate, not an assertion that the
largest table attains it.

Polynomial factors per level add only $O(\lambda\log\lambda)$ to the
time exponent.  Direct leaves cost $2^{O(s/\lambda^{3/4})}$, original
entry routines $2^{O(s^{9/10})}$, and fresh packet routines at most
$\exp(s^{1/21+o(1)})$, all lower order even after the paid recursive
call counts.  Immutable creation-time descriptions prevent hidden
composition of entry routines.  We obtain, under the stated inherited
network interface,
\[
 T(s,\mathcal L),S(s,\mathcal L)
 \le\operatorname{poly}(s+\mathcal L)
       2^{(C_{\rm joint}+o(1))s/\sqrt{\log_2s}}.
\]
The eventual rounded coefficient $31.253$ is certified.  The method uses
subexponential space.  It decreases a leading exponential coefficient,
while leaving the worst-case and fixed-density SAT logarithmic orders
unchanged.
